\documentclass[11pt]{article}
\usepackage[margin=1in]{geometry}
\usepackage{amsmath,amssymb,amsthm,hyperref,enumitem}
\hypersetup{colorlinks=true,urlcolor=blue,linkcolor=blue}
\newtheorem{theorem}{Theorem}
\newtheorem{lemma}[theorem]{Lemma}
\newcommand{\Z}{\mathbb Z}
\title{Uniqueness in the mirror Markov tree}
\author{Principia Math In Progress}
\date{22 September 2026}
\begin{document}
\maketitle
\noindent\textbf{MathDB 373162. Task 1: a proof of Conjecture 6.3 of \cite{mirror}.}
Fix a mirror Markov tree with the initial segment and signed mutations of
\cite[Definitions 4.1--4.2]{mirror}. Its triples are regarded as vertices of
that tree, with the source's ordering by evaluation at $q=1$.
\begin{theorem}\label{main}
Two distinct triples in this tree cannot have the same polynomial as their
maximum-degree entry. In fact, the degree and the value at $q=1$ of that
entry already determine the triple.
\end{theorem}
The proof uses the established fixed-sum ordering theorem for classical
Markov numbers \cite{fixedsum}. It does not assume classical Markov
uniqueness, or positivity of the coefficients of the mirror factors.

\section*{1. Degree on the squared tree}
Index the ordinary Farey tree in $[0,1]$ by reduced fractions $a/b$.
The endpoint labels are $0/1$ and $1/1$, and the first mediant is $1/2$.
Write $m_{a/b}$ for its classical Markov label, with
$m_{0/1}=1$, $m_{1/1}=2$, $m_{1/2}=5$. These are labels at positions:
we do not presume that different positions have different integer labels.
Use the identical indexing for the symmetric squared deformation
$F_{a/b}(q)$ and its mirror factor $f_{a/b}(q)$.
The source establishes
\[
 F_{a/b}(q)=f_{a/b}(q)f_{a/b}(q^{-1}),\qquad
 f_{a/b}(1)=m_{a/b},                                      \tag{1}
\]
and that every $F_{a/b}$ is a symmetric Laurent polynomial with
nonnegative integer coefficients \cite[Theorem 2.2, Section 3.1,
Theorem 4.4 and Notation 4.5]{mirror}.
For a nonzero Laurent polynomial $g$ let $\delta(g)$ be its largest
exponent. For an ordinary polynomial $f\ne0$, its exponent width is
$w(f)=\deg f-\operatorname{val}_q f$. The endpoint terms give
\[
 \delta\bigl(f(q)f(q^{-1})\bigr)=w(f).                    \tag{2}
\]
There is exactly one product contributing to the largest exponent in (2).

\begin{lemma}\label{degree}
For every reduced $a/b\in[0,1]$,
\[
 \delta(F_{a/b})=w(f_{a/b})=a+b-1.                       \tag{3}
\]
At every noninitial squared triple the largest degree is the sum of the
other two degrees plus one, and hence is strictly larger than both.
\end{lemma}
\begin{proof}
The first two triples are $(1,1,1)$ and $(1,1,q+2+q^{-1})$.
The latter has degrees $(0,0,1)$. Suppose a triple $(A,B,C)$,
ordered by its classical labels, has degrees $(r,s,t)$ with
$t=r+s+1$ and $r,s\ge0$. Mutation replacing $A$ gives
\[
 AA'=B^2+(q+q^{-1})BC+C^2.                              \tag{4}
\]
Since $t\ge s+1$, the numerator has largest exponent
$\max\{2s,s+t+1,2t\}=2t$. Its coefficients are nonnegative,
so a tie cannot cancel. Therefore
\[
 \delta(A')=2t-r=s+t+1>t.
\]
Replacing $B$ likewise produces degree $r+t+1>t$.
Induction proves the asserted degree invariant on every branch.
In particular the labels $0/1$, $1/1$, $1/2$ have degrees $0,1,2$.
For a new Farey mediant $(a+c)/(b+d)$, the recurrence just proved is
\[
 \delta(F_{(a+c)/(b+d)})
 =\delta(F_{a/b})+\delta(F_{c/d})+1.
\]
The proposed expression $a+b-1$ satisfies precisely this recurrence and
the endpoint values. Equation (2) finishes the proof.
\end{proof}

\section*{2. The normalization has nonzero constant term}
This step identifies ordinary polynomial degree with the width in (3).
It uses only the source's squared positivity and polynomiality.
The initial mirror polynomials $1$, $1+q$, and $1+2q+2q^2$
have nonzero constant term, as do their reversed choices.
Assume this for the three entries of a triple. Write $d(q)$ for the
entry being replaced, $h(q)$ for the other nonmaximum entry, and
$c(q)$ for the maximum entry. Set
$H=h(q)h(q^{-1})$ and $C=c(q)c(q^{-1})$.
By the induction hypothesis and Lemma~\ref{degree},
\[
 t:=\deg c=\delta(C),\qquad \delta(H)\le t-1.
\]
For either permitted sign $\epsilon\in\{-1,1\}$ the source's
mutation has the form
\[
 p(q)=\frac{q^t\bigl(q^\epsilon H+C\bigr)}{\widetilde d(q)},
 \qquad \widetilde d(q)=q^{\deg d}d(q^{-1}).             \tag{5}
\]
The numerator is an ordinary polynomial: its lowest possible exponents
are $t+\epsilon-\delta(H)\ge0$ and $t-\delta(C)=0$.
Its constant coefficient is positive, since $C$ contributes its positive
endpoint coefficient and $H$ cannot cancel it.
The denominator has nonzero constant coefficient, namely the leading
coefficient of $d$. The source's Theorem 4.4 says that $p$ is a
polynomial. Evaluating the polynomial identity in (5) at zero now gives
$p(0)\ne0$. This proves the induction. Consequently
\[
 \deg f_{a/b}=a+b-1,                                    \tag{6}
\]
and the mediant entry is the unique maximum-degree polynomial of its
triple. The argument applies to every fixed permitted initial choice.

\section*{3. Fixed-sum ordering separates all positions}
We use the following published theorem of Lagisquet, Pelantov\'a,
Tavenas and Vuillon \cite[Theorem 2(3), Proposition 22]{fixedsum}:
if $0<i<a<b$, $\gcd(a,b)=1$, and $\gcd(a-i,b+i)=1$, then
\[
 m_{a/b}<m_{(a-i)/(b+i)}.                               \tag{7}
\]
Thus distinct reduced fractions in $(0,1)$ having the same numerator
plus denominator have different Markov labels. Indeed, given two such
fractions, choose $a$ to be the larger numerator and let $i$ be their
positive numerator difference. The other numerator $a-i$ is positive,
and equality of the sums makes its denominator $b+i$; all hypotheses
of (7) follow. The cited theorem is proved for an extension to
noncoprime pairs as well, so it does not require intermediate fractions
in a comparison chain to be reduced.

Suppose two interior positions give polynomials with equal degree and
equal value at $q=1$. Equation (6) gives equal sums $a+b$; equation (1)
gives equal classical Markov labels. By (7) their reduced fractions
are equal. The ordinary Farey construction has exactly one vertex
whose mediant is any given reduced fraction in $(0,1)$, so the two
triples in the fixed mirror tree coincide.
The exceptional initial triples $(1,1,1)$ and $(1,1,1+q)$ have
maximum degrees $0$ and $1$, whereas every interior triple has
maximum degree at least $2$. They cause no collision with one another
or with an interior triple. This proves Theorem~\ref{main}.

The same reasoning proves injectivity of $a/b\mapsto F_{a/b}(q)$:
equal squared polynomials give equal degrees and
$F_{a/b}(1)=m_{a/b}^2$; positive classical labels then agree.
This is a consequence of the same proof, not a second problem credit.
In particular the proof does not infer the still separate classical
uniqueness assertion by discarding the degree information.

\section*{Verification and scope}
The accompanying exact verifier checks the squared mutation, Farey
degree formula and independent Christoffel-matrix labels on a finite
binary subtree; both initial mirror choices and their signed mutations;
and the fixed-sum order for a bounded set of coprime pairs. These are
supplementary checks. The inductions and the published theorem establish
the unrestricted conclusion. Review is by the authoring agent only.

\section*{Proposed follow-up Tasks 2--5}
These are research directions following the completed Task 1.
\begin{enumerate}[label=\textbf{Task \arabic*.},start=2,leftmargin=*]
\item Construct an algorithm recovering the Farey fraction from
$(\deg f,f(1))$ by the fixed-sum ordering, and give bit-complexity bounds.
\item Determine whether the degree and two endpoint coefficients already
identify every mirror polynomial in a fixed tree.
\item Derive a closed Christoffel-word formula for the leading and constant
coefficients of $f_{a/b}$, including their dependence on the initial choice.
\item For the squared recurrence with $q+q^{-1}$ replaced by a symmetric
Laurent polynomial, characterize the substitutions for which a degree
statistic and an integer specialization still separate all Farey positions.
\end{enumerate}

\begin{thebibliography}{9}
\bibitem{mirror} L. Bittmann, P. Jouteur, E. Kantarc{\i} O\u{g}uz,
M. Molander and E. Y{\i}ld{\i}r{\i}m,
\emph{A mirror deformation of Markov numbers},
\href{https://arxiv.org/abs/2602.14802}{arXiv:2602.14802v1}, 2026.
\bibitem{fixedsum} C. Lagisquet, E. Pelantov\'a, S. Tavenas and L. Vuillon,
\emph{On the Markov numbers: fixed numerator, denominator, and sum conjectures},
Advances in Applied Mathematics \textbf{130} (2021), 102227,
\href{https://doi.org/10.1016/j.aam.2021.102227}{doi:10.1016/j.aam.2021.102227}.
Theorem 2 and Proposition 22 in the complete preprint
\href{https://arxiv.org/abs/2010.10335}{arXiv:2010.10335v1}.
\end{thebibliography}
\end{document}
